\chapter{Databases, SQL Concurrency, Storage Internals, and ORMs}
\label{chap:databases}

Imagine a hospital scheduling table with two rows:

\begin{center}
\small\ttfamily
\begin{tabular}{@{}lll@{}}
\toprule
\textbf{doctor\_id} & \textbf{name} & \textbf{on\_call} \\ \midrule
1 & Dr.\ Alice & true \\
2 & Dr.\ Bob & true \\ \bottomrule
\end{tabular}
\end{center}

Hospital policy says at least one doctor must stay on call at all times. Both Alice and Bob try to go off-duty at the exact same millisecond: each browser runs \texttt{SELECT count(*) FROM doctors WHERE on\_call = true;}, sees \texttt{2}, decides ``it is safe for me to step off'', and runs \texttt{UPDATE doctors SET on\_call = false WHERE doctor\_id = ...;}. When both queries commit, \textbf{zero doctors are on call}---even though neither transaction saw uncommitted data and both queries looked completely valid!

To understand why concurrency bugs like that happen---and how databases keep data consistent without grinding to a halt---we work from the logical table level down into PostgreSQL's B-Trees, 8\,KB heap pages, Write-Ahead Log, and ORM boundaries.

% ==========================================
\section{ACID, Concurrency Anomalies, and Isolation Levels}

A database \textbf{transaction} groups SQL statements into one logical unit governed by \textbf{ACID}:
\begin{itemize}[leftmargin=15pt]
  \item \textbf{Atomicity:} Either every write in the transaction succeeds (\texttt{COMMIT}), or every write is undone (\texttt{ROLLBACK}).
  \item \textbf{Consistency:} Constraints (\texttt{CHECK}, foreign keys, unique indexes) hold before and after every transaction.
  \item \textbf{Isolation:} Concurrent transactions do not corrupt each other's view of data.
  \item \textbf{Durability:} Once \texttt{COMMIT} returns success, the change survives a power failure.
\end{itemize}

\subsection{Four Concurrency Anomalies Defined by Timelines}

Stronger isolation prevents more anomalies at higher coordination cost. Let's define all four anomalies with concrete two-transaction ($T_1, T_2$) interleavings:

\begin{enumerate}[leftmargin=15pt]
  \item \textbf{Dirty Read:} $T_1$ writes \texttt{balance = 500} (uncommitted) $\rightarrow$ $T_2$ reads \texttt{500} $\rightarrow$ $T_1$ executes \texttt{ROLLBACK} back to \texttt{1000}. $T_2$ acted on money that never existed.
  \item \textbf{Non-Repeatable Read:} $T_1$ reads \texttt{balance = 1000} $\rightarrow$ $T_2$ updates \texttt{balance = 500} and \texttt{COMMIT}s $\rightarrow$ $T_1$ reads the same row again and sees \texttt{500} inside the same transaction.
  \item \textbf{Phantom Read:} $T_1$ runs \texttt{SELECT count(*) FROM orders WHERE status = 'NEW'} (sees 3) $\rightarrow$ $T_2$ \texttt{INSERT}s a 4th \texttt{'NEW'} order and \texttt{COMMIT}s $\rightarrow$ $T_1$ re-runs the range query and sees 4 rows.
  \item \textbf{Write Skew:} Our hospital on-call bug above! $T_1$ reads Alice + Bob both \texttt{on\_call=true}, $T_2$ reads Alice + Bob both \texttt{on\_call=true}; $T_1$ updates row 1 while $T_2$ updates row 2. Because they modify \emph{disjoint rows} based on a shared premise, both commit and break the invariant!
\end{enumerate}

\begin{table}[htbp]
\centering
\small\sffamily
\begin{tabular}{@{}lllll@{}}
\toprule
\textbf{PostgreSQL Level} & \textbf{Dirty Read} & \textbf{Non-Repeatable Read} & \textbf{Phantom Read} & \textbf{Write Skew} \\ \midrule
\textbf{Read Committed} (Default) & Prevented & Allowed & Allowed & Allowed \\
\textbf{Repeatable Read} & Prevented & Prevented & Prevented (PG) & Allowed \\
\textbf{Serializable} (SSI) & Prevented & Prevented & Prevented & Prevented \\ \bottomrule
\end{tabular}
\caption{PostgreSQL isolation levels and prevented anomalies.}
\label{tab:v1-ch05-isolation}
\end{table}

\subsection{2PL vs.\ MVCC, and \texttt{SELECT FOR UPDATE SKIP LOCKED}}

Databases enforce isolation using two contrasting philosophies:
\begin{itemize}[leftmargin=15pt]
  \item \textbf{Two-Phase Locking (2PL):} Transactions acquire shared locks to read and exclusive locks to write, holding all locks until \texttt{COMMIT}/\texttt{ROLLBACK}. Readers block writers and writers block readers---simple to reason about, but slow under read-heavy traffic (studied further in distributed transactions in Volume~IV).
  \item \textbf{Multi-Version Concurrency Control (MVCC):} Used by PostgreSQL. Instead of overwriting rows in place, an \texttt{UPDATE} writes a \emph{new version} of the row stamped with transaction IDs, so \textbf{readers never block writers, and writers never block readers} \cite{kleppmann2017designing}.
\end{itemize}

\begin{intuitionbox}[Library Ledger (2PL) vs.\ Timestamped Photocopies (MVCC)]
Under 2PL, if someone is editing a page in the master ledger, nobody else can even read that page until they finish. Under MVCC, readers receive a consistent snapshot photocopy of the ledger as of the moment their query or transaction began, while writers append new dated versions alongside it.
\end{intuitionbox}

When building job queues or flash-sale reservations where workers \emph{must} claim disjoint rows exclusively, explicit row locks combine MVCC speed with safe exclusion:

\begin{lstlisting}[language=SQL,caption={Lock-Contention-Free Queue Worker Pattern with \texttt{SKIP LOCKED}}]
BEGIN;
SELECT id, payload
FROM job_queue
WHERE status = 'PENDING'
ORDER BY priority DESC, id ASC
LIMIT 1
FOR UPDATE SKIP LOCKED;

UPDATE job_queue SET status = 'IN_PROGRESS' WHERE id = :claimed_id;
COMMIT;
\end{lstlisting}

\textbf{Concrete 2-worker trace:} Suppose rows 101 and 102 are \texttt{'PENDING'}. Worker 1 runs the \texttt{SELECT ... FOR UPDATE SKIP LOCKED}, locks row 101, and starts processing. Milliseconds later Worker 2 runs the same query: instead of blocking behind Worker 1's lock on row 101 (as plain \texttt{FOR UPDATE} would), \texttt{SKIP LOCKED} skips row 101 immediately, locks row 102, and returns it to Worker 2 with zero lock waiting!

% ==========================================
\section{B-Tree Indexes and Query Plan Profiling (\texttt{EXPLAIN})}

Scanning millions of rows sequentially from disk for one user ID wastes I/O. PostgreSQL accelerates lookups with \textbf{Lehman and Yao B-Tree} indexes: balanced trees of 8\,KB index pages consisting of a \textbf{Metapage} (block 0 pointing to the root), \textbf{Internal Branch Pages} (routing keys + \textbf{High Key} upper bounds), and doubly-linked \textbf{Leaf Pages} holding indexed values paired with physical heap row pointers (\textbf{TIDs}).

\begin{figure}[htbp]
\centering
\begin{tikzpicture}[
  scale=0.88, transform shape,
  node distance=0.8cm and 1.1cm,
  node/.style={rectangle, draw=headingblue, fill=codebg, rounded corners=2pt, thick, align=center, inner sep=5pt, font=\small\sffamily},
  leaf/.style={rectangle, draw=intuitionTeal, fill=intuitionBg, rounded corners=2pt, thick, align=center, inner sep=5pt, font=\small\sffamily\bfseries},
  line/.style={draw=headingblue, -{Stealth[length=2.3mm]}, thick}
]
  \node[node] (root) {\textbf{Root Page (Block 3)}\\\small Keys: [50, 100]};
  \node[node, below=0.9cm of root, xshift=-3.4cm] (b1) {\textbf{Branch Page}\\\small Keys: [10, 30]};
  \node[node, below=0.9cm of root, xshift=3.4cm] (b2) {\textbf{Branch Page}\\\small Keys: [60, 80]};
  
  \node[leaf, below=0.9cm of b1, xshift=-1.4cm] (l1) {Leaf [1..9]\\TIDs};
  \node[leaf, below=0.9cm of b1, xshift=1.4cm] (l2) {Leaf [10..29]\\TIDs};
  \node[leaf, below=0.9cm of b2, xshift=-1.4cm] (l3) {Leaf [50..59]\\TIDs};
  \node[leaf, below=0.9cm of b2, xshift=1.4cm] (l4) {Leaf [100..]\\TIDs};

  \draw[line] (root) -- (b1);
  \draw[line] (root) -- (b2);
  \draw[line] (b1) -- (l1);
  \draw[line] (b1) -- (l2);
  \draw[line] (b2) -- (l3);
  \draw[line] (b2) -- (l4);
  \draw[line, <->, intuitionTeal] (l1) -- (l2);
  \draw[line, <->, intuitionTeal] (l2) -- (l3);
  \draw[line, <->, intuitionTeal] (l3) -- (l4);
\end{tikzpicture}
\caption{B-Tree page hierarchy and doubly-linked leaf pages enabling ordered range scans.}
\label{fig:v1-ch05-btree-arch}
\end{figure}

\begin{table}[htbp]
\centering
\small\sffamily
\begin{tabular}{@{}lll@{}}
\toprule
\textbf{Access Method} & \textbf{Mechanics} & \textbf{Best When} \\ \midrule
\textbf{Sequential Scan} & Reads every 8\,KB heap page linearly. & Query matches $>15\%$ of rows or tiny table. \\
\textbf{Index Scan} & Walks B-Tree $\to$ fetches each TID from heap. & Selective lookups (\texttt{id = 42}). \\
\textbf{Bitmap Index Scan} & Collects matching TIDs, sorts by block, reads heap sequentially. & Moderate ranges ($2\%\text{--}10\%$) avoiding random disk seeks. \\
\textbf{Index-Only Scan} & Answers purely from index if Visibility Map bit $= 1$. & Covering indexes where all needed columns are indexed. \\ \bottomrule
\end{tabular}
\caption{PostgreSQL table and index scan strategies.}
\label{tab:v1-ch05-scans}
\end{table}

\subsection{Reading \texttt{EXPLAIN (ANALYZE, BUFFERS)} Bottom-Up}

Always read PostgreSQL execution plans \textbf{from the innermost indented child nodes upward} toward the root, because child operators feed rows up to parent operators and buffer counts are cumulative:

\begin{lstlisting}[language=SQL,caption={Query Under Analysis}]
EXPLAIN (ANALYZE, BUFFERS)
SELECT u.name, COUNT(o.id)
FROM users u
JOIN orders o ON o.user_id = u.id
WHERE o.created_at >= '2026-01-01'
GROUP BY u.name;
\end{lstlisting}

\begin{lstlisting}[language=sh,caption={Execution Plan Output (Read Child Nodes Bottom-Up)}]
HashAggregate  (cost=1450.20..1520.40 rows=7020 width=40) (actual time=12.450..14.120 rows=6850 loops=1)
  Group Key: u.name
  Buffers: shared hit=380 read=12
  ->  Hash Join  (cost=210.00..1380.00 rows=14040 width=36) (actual time=2.100..9.500 rows=14200 loops=1)
        Hash Cond: (o.user_id = u.id)
        Buffers: shared hit=380 read=12
        ->  Seq Scan on orders o  (cost=0.00..1020.00 rows=14040 width=8) (actual time=0.040..4.100 rows=14200 loops=1)
              Filter: (created_at >= '2026-01-01'::timestamp)
              Buffers: shared hit=290 read=10
        ->  Hash  (cost=120.00..120.00 rows=7200 width=36) (actual time=1.950..1.950 rows=7200 loops=1)
              Buckets: 8192  Batches: 1  Memory Usage: 480kB
              Buffers: shared hit=90 read=2
              ->  Seq Scan on users u  (cost=0.00..120.00 rows=7200 width=36) (actual time=0.020..1.100 rows=7200 loops=1)
                    Buffers: shared hit=90 read=2
Planning Time: 0.280 ms
Execution Time: 15.200 ms
\end{lstlisting}

\textbf{Bottom-up trace:}
\begin{enumerate}[leftmargin=15pt]
  \item Innermost child \texttt{Seq Scan on users u} reads $90\text{ (hit)} + 2\text{ (read)} = 92$ 8\,KB pages and emits 7,200 rows into \texttt{Hash} ($480\text{\,kB}$ in memory, still \texttt{hit=90 read=2}).
  \item Sibling \texttt{Seq Scan on orders o} reads $290 + 10 = 300$ pages and emits 14,200 matching rows.
  \item \texttt{Hash Join} combines both inputs, so its cumulative buffer counts equal $(90+290)=380$ hits and $(2+10)=12$ reads.
  \item Finally, \texttt{HashAggregate} groups the 14,200 joined rows in memory into 6,850 users without reading extra table pages (\texttt{shared hit=380 read=12}). Notice planner estimates (\texttt{rows=14040}) closely match \texttt{actual rows=14200}; huge mismatches signal stale statistics needing \texttt{ANALYZE}.
\end{enumerate}

% ==========================================
\section{PostgreSQL Physical Storage: 8\,KB Pages, Bloat, and Redo WAL}

\subsection{Inside an 8\,KB Heap Page and \texttt{HeapTupleHeaderData}}

PostgreSQL splits each table file into fixed \textbf{8,192-byte (8\,KB) pages}:

\begin{figure}[htbp]
\centering
\begin{tikzpicture}[
  scale=0.88, transform shape,
  node distance=0.08cm,
  blk/.style={rectangle, draw=headingblue, fill=codebg, thick, minimum height=0.95cm, align=center, font=\scriptsize\sffamily}
]
  \node[blk, minimum width=2.4cm, fill=intuitionBg] (hdr) {\textbf{PageHeaderData}\\(24 bytes)};
  \node[blk, minimum width=2.5cm, right=of hdr] (ptrs) {\textbf{Line Pointers}\\(4 bytes each)};
  \node[blk, minimum width=3.5cm, right=of ptrs, fill=gray!10] (free) {\textbf{Free Space}\\\small Shrinks inward};
  \node[blk, minimum width=2.1cm, right=of free] (t2) {\textbf{Tuple 2}};
  \node[blk, minimum width=2.1cm, right=of t2] (t1) {\textbf{Tuple 1}};

  \draw[-{Stealth[length=2mm]}, thick, intuitionTeal] (ptrs.east) -- ++(0.5,0) node[above, font=\tiny\sffamily] {Grows down};
  \draw[-{Stealth[length=2mm]}, thick, warningRed] (t2.west) -- ++(-0.5,0) node[above, font=\tiny\sffamily] {Grows up};
  \node[draw=headingblue, dashed, fit=(hdr) (ptrs) (free) (t2) (t1), inner sep=5pt] {};
\end{tikzpicture}
\caption{PostgreSQL 8\,KB heap page layout: 24-byte header, downward line pointers, upward tuples.}
\label{fig:v1-ch05-postgres-page}
\end{figure}

Every row is addressed physically by a \textbf{Tuple Identifier (TID)} $= (\text{block\_number}, \text{line\_pointer\_index})$ and carries a \textbf{23-byte \texttt{HeapTupleHeaderData}} (plus a 1+-byte NULL bitmap padded to 8-byte alignment, making the header 24 bytes for up to 8 columns). Inside that header sit \textbf{\texttt{t\_xmin}} (the XID that inserted this row version) and \textbf{\texttt{t\_xmax}} (the XID that deleted or updated it, or $0$ if live).

\subsection{Dead Tuples, Autovacuum, and XID Wraparound}

Because an \texttt{UPDATE} leaves the old version on disk with \texttt{t\_xmax} set and appends a brand-new tuple, frequent updates create \textbf{dead tuples} (\textbf{table bloat}). The \textbf{Autovacuum} worker reclaims dead space, updates the \textbf{Free Space Map (FSM)} and \textbf{Visibility Map (VM)}, and \textbf{freezes} ancient tuples (\texttt{HEAP\_XMIN\_FROZEN}) before PostgreSQL's 32-bit unsigned transaction counter ($\approx 4.29$ billion) wraps around modulo $2^{32}$---which would otherwise force PostgreSQL into an emergency read-only shutdown!

\subsection{Shared Buffers, Clock Sweep, and Redo-Only WAL Recovery}

PostgreSQL caches 8\,KB pages in \textbf{Shared Buffers} (evicted via a ring-hand \textbf{Clock Sweep} usage counter $0\dots 5$). Before modifying any page in memory, PostgreSQL writes a sequential record to the \textbf{Write-Ahead Log (WAL)} and calls \texttt{fsync()} on \texttt{COMMIT} \cite{petrov2019database}.

\begin{productionwarningbox}[Why PostgreSQL Crash Recovery Performs Redo Only (Unlike ARIES Undo)]
Traditional database engines that overwrite rows in place (such as MySQL InnoDB or IBM DB2) use the three-phase \textbf{ARIES} recovery algorithm: *Analysis* $\to$ *Redo* $\to$ *Undo* (to roll back uncommitted in-place edits).

\textbf{PostgreSQL does NOT run an Undo pass during crash recovery.} Because MVCC never overwrites old tuples in place, an uncommitted transaction's half-written tuples simply have an \texttt{xmin} not marked committed in the commit log (\texttt{pg\_xact}). After a power failure, PostgreSQL replays WAL records forward from the last checkpoint's REDO point (\textbf{Redo pass only}); uncommitted tuples are automatically invisible to every future query and are reclaimed later by \texttt{VACUUM}.
\end{productionwarningbox}

\begin{realworldbox}[Diagnosing Flash-Sale Bloat, Missing Index, and Lock Contention]
\textbf{Symptom:} During a product launch, PostgreSQL CPU hits 98\%, update latency climbs to 4 seconds, and table disk size grows from 4\,GB to 62\,GB in one afternoon.

\textbf{Diagnostic Steps:}
\begin{enumerate}[leftmargin=15pt,itemsep=2pt]
  \item Query \texttt{pg\_stat\_user\_tables}: see \texttt{n\_dead\_tup = 48{,}000{,}000} and \texttt{last\_autovacuum = NULL} because a long-running analytics query holds an old snapshot open, blocking autovacuum from reclaiming dead rows.
  \item Query \texttt{pg\_stat\_activity} for \texttt{wait\_event\_type = 'Lock'}: dozens of checkout connections pile up waiting on the same SKU row lock, or run full sequential scans on a JSONB filter lacking an expression/GIN index.
  \item Fix: Terminate stale reporting snapshots, tune \texttt{autovacuum\_vacuum\_scale\_factor = 0.05}, add targeted indexes, and use atomic conditional updates or \texttt{SKIP LOCKED}.
\end{enumerate}
\end{realworldbox}

% ==========================================
\section{ORM Design and Multi-Tenant Row-Level Security}

\subsection{Active Record vs.\ Data Mapper and the N+1 Problem}

\begin{figure}[htbp]
\centering
\begin{tikzpicture}[
  scale=0.88, transform shape,
  node distance=0.7cm and 1.4cm,
  box/.style={rectangle, draw=headingblue, fill=codebg, rounded corners=2pt, thick, align=center, inner sep=5pt, font=\small\sffamily},
  accent/.style={rectangle, draw=intuitionTeal, fill=intuitionBg, rounded corners=2pt, thick, align=center, inner sep=5pt, font=\small\sffamily\bfseries},
  line/.style={draw=headingblue, -{Stealth[length=2.3mm]}, thick}
]
  \node[accent, minimum width=5cm] (ar) {\textbf{Active Record (Eloquent, Django)}\\Model holds data + SQL save/query};
  \node[box, below=0.45cm of ar, minimum width=5cm] (ardb) {Table \texttt{orders}};
  \draw[line, <->] (ar) -- node[right, font=\tiny\sffamily] {\texttt{\$order->save()}} (ardb);

  \node[accent, right=1.6cm of ar, minimum width=5.4cm] (dm_ent) {\textbf{Data Mapper Entity (Doctrine, SQLAlchemy)}\\Pure domain object (no DB methods)};
  \node[box, below=0.45cm of dm_ent, minimum width=5.4cm] (dm_map) {\textbf{Unit of Work / Session}\\Tracks dirty objects, batches SQL};
  \node[box, below=0.45cm of dm_map, minimum width=5.4cm] (dm_db) {Table \texttt{orders}};
  \draw[line] (dm_ent) -- (dm_map);
  \draw[line, <->] (dm_map) -- (dm_db);
\end{tikzpicture}
\caption{Active Record vs.\ Data Mapper persistence architectures.}
\label{fig:v1-ch05-orm-patterns}
\end{figure}

Both ORM families easily trigger the \textbf{N+1 query problem}: fetching 100 orders (1 query) and reading \texttt{\$order->customer->name} inside a loop fires 100 extra \texttt{SELECT} round trips. Fix it with \textbf{eager loading} (\texttt{WHERE id IN (...)}, turning 101 queries into 2) or batching across async requests via the \textbf{DataLoader} pattern.

\subsection{Row-Level Security (RLS) and \texttt{SET LOCAL} with PgBouncer}

In multi-tenant backends, PostgreSQL \textbf{Row-Level Security (RLS)} acts as a database-enforced safety net: even if an ORM query forgets \texttt{WHERE tenant\_id = ...}, PostgreSQL automatically filters every row:

\begin{lstlisting}[language=SQL,caption={Enforcing Tenant Isolation Safely with RLS and \texttt{SET LOCAL}}]
ALTER TABLE invoices ENABLE ROW LEVEL SECURITY;
CREATE POLICY tenant_isolation ON invoices
    USING (tenant_id = current_setting('app.current_tenant', true)::uuid);

-- Inside each request transaction:
BEGIN;
SET LOCAL app.current_tenant = '11111111-2222-3333-4444-555555555555';
SELECT * FROM invoices;  -- Automatically filtered to that tenant only
COMMIT;                  -- SET LOCAL resets automatically when the transaction ends!
\end{lstlisting}

\begin{productionwarningbox}[Why \texttt{SET} Leaks Tenants Behind PgBouncer Transaction Pooling]
Connection poolers like \textbf{PgBouncer} in \texttt{transaction} pooling mode hand the exact same PostgreSQL backend socket to Tenant A for Transaction 1 and Tenant B for Transaction 2. Using session-scoped \texttt{SET app.current\_tenant = ...} leaves Tenant A's ID sticky on the socket for Tenant B! Always use \textbf{\texttt{SET LOCAL}} inside an explicit \texttt{BEGIN ... COMMIT} block so PostgreSQL wipes the setting automatically at transaction end (explored further in Volume~V).
\end{productionwarningbox}

\begin{tryityourselfbox}[Trace MVCC Visibility Across Two Transactions]
Row \texttt{id=1} has \texttt{(xmin=100, xmax=0, balance=500)}. Transaction \texttt{XID=105} begins at \texttt{REPEATABLE READ} and runs \texttt{SELECT balance FROM accounts WHERE id=1}. Meanwhile, Transaction \texttt{XID=106} updates \texttt{balance = 300} (creating version \texttt{xmin=106, xmax=0} and setting \texttt{xmax=106} on the old version) and commits. Now \texttt{XID=105} runs \texttt{SELECT balance} a second time. Which tuple version does \texttt{XID=105} see, and what changes if \texttt{XID=105} was \texttt{READ COMMITTED}?
\end{tryityourselfbox}

% ==========================================
\section{Chapter Summary \protect\& Key Takeaways}

\begin{itemize}[leftmargin=15pt]
  \item \textbf{Anomalies \& Isolation:} Dirty Reads, Non-Repeatable Reads, Phantoms, and Write Skew define isolation guarantees. PostgreSQL MVCC prevents dirty/non-repeatable/phantom reads at \texttt{Repeatable Read} and requires \texttt{Serializable} (SSI) for Write Skew.
  \item \textbf{Locking \& Plans:} \texttt{SELECT ... FOR UPDATE SKIP LOCKED} lets concurrent queue workers claim disjoint jobs without waiting; \texttt{EXPLAIN (ANALYZE, BUFFERS)} must be read bottom-up with cumulative buffer hits/reads.
  \item \textbf{Storage \& Recovery:} 8\,KB pages pack 24-byte headers, 4-byte line pointers, and 8-byte-aligned tuples with \texttt{xmin}/\texttt{xmax}. Autovacuum clears dead tuples and prevents 32-bit XID wraparound. PostgreSQL WAL crash recovery runs a \textbf{Redo pass only} (no ARIES Undo pass).
  \item \textbf{ORMs \& Multi-Tenancy:} Prevent N+1 loops with eager \texttt{WHERE IN} queries, and pair PostgreSQL RLS policies with transaction-scoped \texttt{SET LOCAL} behind connection poolers.
\end{itemize}

% ==========================================
\section{Exercises and Real-World Scenarios}

\subsection{Mechanics Exercises}
\begin{enumerate}[leftmargin=15pt]
  \item \textbf{Exercise~\thechapter.1:} Calculate the exact physical heap storage required by 1,000,000 uncompressed rows in a table defined as \texttt{(id INT, created\_at TIMESTAMP, label VARCHAR(32), active BOOLEAN)}, assuming every \texttt{label} is 32 ASCII characters (1-byte short varlena header + 32 bytes = 33 bytes), taking into account PostgreSQL's 24-byte page header, 4-byte line pointers, 24-byte tuple header (23B + 1B NULL bitmap), and 8-byte \texttt{MAXALIGN} alignment padding.
  \item \textbf{Exercise~\thechapter.2:} Explain the exact difference between an Index Scan and an Index-Only Scan, and name the exact data structure bit that forces an Index-Only Scan to fall back to reading the heap page.
\end{enumerate}

\subsection{Debugging Exercise}
\begin{enumerate}[leftmargin=15pt]
  \item \textbf{Exercise~\thechapter.3:} This query runs a slow Sequential Scan across 50,000,000 rows even though an index exists:
\begin{lstlisting}[language=SQL]
CREATE INDEX idx_audit_logs_metadata ON audit_logs (metadata);
SELECT * FROM audit_logs WHERE metadata->>'tenant_id' = 'acme_corp';
\end{lstlisting}
  Explain why PostgreSQL ignores \texttt{idx\_audit\_logs\_metadata} and provide two valid indexing fixes.
\end{enumerate}

\subsection{Architecture Design Problem}
\begin{enumerate}[leftmargin=15pt]
  \item \textbf{Exercise~\thechapter.4:} Design a flash-sale inventory reservation schema and transaction flow in PostgreSQL for 5,000 concurrent shoppers competing for 500 units, guaranteeing zero overselling, zero deadlocks, and minimal lock hold times.
\end{enumerate}

% ==========================================
\section{Detailed Solutions}

\subsection{Solution to Try It Yourself: MVCC \texttt{xmin}/\texttt{xmax} Trace}
\begin{itemize}[leftmargin=15pt]
  \item Under \textbf{\texttt{REPEATABLE READ}}, \texttt{XID=105} takes one transaction-level snapshot at its first query (when \texttt{XID=106} had not committed yet). On its second \texttt{SELECT}, it sees the old tuple (\texttt{xmin=100, xmax=106}) as still live because \texttt{106} is not in its snapshot, and ignores the new tuple (\texttt{xmin=106}). Result: \textbf{500}.
  \item Under \textbf{\texttt{READ COMMITTED}}, \texttt{XID=105} takes a fresh snapshot at the start of \emph{each statement}. Since \texttt{106} is committed before the second \texttt{SELECT}, it sees the old tuple as deleted and returns the new tuple: \textbf{300}.
\end{itemize}

\subsection{Solution~\thechapter.1: Exact 8-Byte Aligned PostgreSQL Heap Storage Math}
\begin{enumerate}[leftmargin=15pt]
  \item \textbf{Tuple Header + NULL Bitmap:} \texttt{HeapTupleHeaderData} is 23 bytes + 1 byte NULL bitmap (for 4 columns) $= \mathbf{24\text{ bytes}}$ (already a multiple of 8).
  \item \textbf{Column Layout \& Alignment Padding (column order: \texttt{id, created\_at, label, active}):}
  \begin{itemize}
    \item \texttt{id} (\texttt{INT}, 4 bytes, 4-byte aligned): offsets $24\dots 27$ (4 bytes).
    \item \texttt{created\_at} (\texttt{TIMESTAMP}, 8 bytes, 8-byte aligned): inserts \textbf{4 bytes alignment padding} at $28\dots 31$, then occupies $32\dots 39$ (8 bytes).
    \item \texttt{label} (\texttt{VARCHAR(32)}, 1-byte header + 32 bytes = 33 bytes, 1-byte aligned): offsets $40\dots 72$ (33 bytes).
    \item \texttt{active} (\texttt{BOOLEAN}, 1 byte, 1-byte aligned): offset $73$ (1 byte).
  \end{itemize}
  \item \textbf{Total Tuple \& Page Slot Size:} Unaligned tuple end $= 74$ bytes. PostgreSQL pads every tuple length to 8-byte \texttt{MAXALIGN}: $\lceil 74 / 8 \rceil \times 8 = \mathbf{80\text{ bytes}}$. Adding the 4-byte \texttt{ItemIdData} line pointer gives $\mathbf{84\text{ bytes per row}}$ on the page. (Reordering columns largest-fixed-first---putting \texttt{created\_at} before \texttt{id}---eliminates the 4-byte padding gap, shrinking each tuple to 72 bytes + 4-byte pointer = 76 bytes!)
  \item \textbf{Page Capacity \& Total Heap Size:}
  \begin{itemize}
    \item Usable space per page: $8192 - 24 = 8168\text{ bytes}$.
    \item Tuples per page: $\lfloor 8168 / 84 \rfloor = \mathbf{97\text{ tuples/page}}$.
    \item Total pages for $1{,}000{,}000$ rows: $\lceil 1{,}000{,}000 / 97 \rceil = \mathbf{10{,}310\text{ pages}}$.
    \item Total heap size: $10{,}310 \times 8192\text{ bytes} = 84{,}459{,}520\text{ bytes} \approx \mathbf{80.55\text{ MiB}}$ ($\approx 84.46\text{\,MB}$ decimal).
  \end{itemize}
\end{enumerate}

\subsection{Solution~\thechapter.2: Index Scan vs.\ Index-Only Scan}
An Index Scan walks the B-Tree and visits the heap page for every matching TID to read missing columns and check \texttt{xmin}/\texttt{xmax} visibility. An Index-Only Scan returns columns straight from the B-Tree leaf pages, checking the \textbf{Visibility Map (VM)} bit for each heap page: if the VM all-visible bit is \texttt{1}, the heap read is skipped; if \texttt{0} (recent uncommitted or unvacuumed changes on that page), PostgreSQL falls back to reading the heap page to verify MVCC visibility.

\subsection{Solution~\thechapter.3: Fixing JSONB Expression Indexing}
A standard B-Tree on \texttt{(metadata)} orders whole JSONB objects, not the extracted key \texttt{metadata->>'tenant\_id'}.
\begin{lstlisting}[language=SQL]
-- Fix A: B-Tree expression index matching the exact query predicate
CREATE INDEX idx_audit_logs_tenant ON audit_logs ((metadata->>'tenant_id'));

-- Fix B: GIN index + containment operator (@>) for flexible JSONB keys
CREATE INDEX idx_audit_logs_gin ON audit_logs USING GIN (metadata jsonb_path_ops);
-- Query using: WHERE metadata @> '{"tenant_id": "acme_corp"}'
\end{lstlisting}

\subsection{Solution~\thechapter.4: Deadlock-Free Flash-Sale Reservation Flow}
\begin{lstlisting}[language=SQL]
CREATE TABLE inventory_items (
    sku VARCHAR(64) PRIMARY KEY,
    total_stock INT NOT NULL CHECK (total_stock >= 0),
    reserved_stock INT NOT NULL DEFAULT 0 CHECK (reserved_stock <= total_stock)
);

CREATE TABLE reservations (
    id BIGSERIAL PRIMARY KEY,
    sku VARCHAR(64) NOT NULL REFERENCES inventory_items(sku),
    user_id BIGINT NOT NULL,
    created_at TIMESTAMPTZ NOT NULL DEFAULT now()
);

-- Single-roundtrip atomic claim in default READ COMMITTED:
BEGIN;
WITH claimed AS (
    UPDATE inventory_items
    SET reserved_stock = reserved_stock + 1
    WHERE sku = 'CONSOLE-X' AND (total_stock - reserved_stock) >= 1
    RETURNING sku
)
INSERT INTO reservations (sku, user_id)
SELECT sku, :user_id FROM claimed
RETURNING id;
COMMIT;
\end{lstlisting}
\textbf{Why this prevents overselling and deadlocks:} PostgreSQL takes a row-level lock on the single \texttt{CONSOLE-X} row during the atomic \texttt{UPDATE ... WHERE (total\_stock - reserved\_stock) >= 1}. Under \texttt{READ COMMITTED}, waiting transactions automatically re-evaluate the \texttt{WHERE} predicate against the latest committed \texttt{reserved\_stock} once the lock releases. Once \texttt{reserved\_stock == 500}, the \texttt{UPDATE} matches 0 rows, \texttt{claimed} is empty, no reservation row is inserted, and the API returns \texttt{409 Conflict} with zero multi-row deadlock cycles.

\begin{labbridgebox}[Practice Real SQL Transactions \& Async ORMs]
\begin{itemize}[leftmargin=15pt,itemsep=2pt]
  \item \texttt{exercises/volume1/ex02\_php\_pdo\_prepared\_transactions.php} \& \texttt{ex04\_sql\_mvcc\_isolation.sql}: Practice PDO prepared statements, transaction rollbacks, and MVCC row locking.
  \item \texttt{projects/p01-php-db-integration/}: Full PHP 8.3 order placement service with isolation levels and row locking.
  \item \texttt{projects/p04-async-db-sqlalchemy/}: Async SQLAlchemy 2.0 + \texttt{asyncpg} connection pooling and batched upserts (\texttt{ON CONFLICT DO UPDATE}).
\end{itemize}
\end{labbridgebox}

\section*{Further Reading}
See Kleppmann on isolation anomalies and MVCC \cite{kleppmann2017designing} and Petrov's \textit{Database Internals} on B-Trees and WAL \cite{petrov2019database}.
